% Transcription — fonds Alexandre Grothendieck, shelfmark 161-2, batch 3 (pages 41-60).
% Established from the facsimile archives/batches/161-2/batch-03.pdf (a mirror of
% https://grothendieck.umontpellier.fr/161-2.pdf), per the skill
% transcribe-grothendieck. Page 53 is blank and page 57 carries a single scrawl
% with no mathematics — both absent. The batch holds three runs: the spiral
% notebook on the classifying topoi of ring axioms, continued from batch 2
% (pp. 41-50, items 11 to 22 of its list); two loose lined sheets on flatness
% and module structures over a topos (pp. 51-52); and a run of yellowed sheets on
% accessible objects and cardinal filtrations, keyed to SGA 4 exposé I § 9
% (pp. 54-56, 58-60).
% Pass: Fable 5.1 (claude-fable-5-1), 2026-09-04 — first pass, unchecked against the pages by a human.
\documentclass[11pt,a4paper]{article}
% Le préambule commun à toutes les transcriptions du fonds Grothendieck.
%
% Il tient en une idée : **l'incertitude se compose, elle ne se résout pas.**
% Une transcription qui lisse un mot douteux a détruit la seule chose pour
% laquelle on la faisait — un lecteur doit pouvoir savoir, à chaque endroit,
% ce qui était sur le feuillet et ce qui a été ajouté. Les six macros qui
% suivent sont donc l'appareil critique, et non de la décoration.
%
% Les mêmes commandes sont reconnues par `scripts/render.mjs`, qui produit la
% vue de lecture du site. Une macro ajoutée ici doit l'être aussi là-bas,
% faute de quoi le rendu échouera bruyamment — ce qui est voulu : un rendu
% silencieusement faux est pire qu'un rendu absent.

\ProvidesPackage{grothendieck}[2026/08/08 Transcriptions du fonds Grothendieck]

\usepackage{amsmath,amssymb,amsthm}
\usepackage{mathrsfs}
\usepackage{stmaryrd}
% Les diagrammes commutatifs sont partout dans ce fonds : c'est la langue
% même de la géométrie algébrique de Grothendieck, et une transcription qui
% les rendrait en prose ne transcrirait rien.
\usepackage{tikz-cd}
% Français par défaut — c'est la langue des feuillets — mais l'anglais est
% chargé aussi : la traduction et le résumé sont des documents anglais, et la
% typographie française (espaces avant les deux-points) y serait une faute.
% Ces documents basculent par \selectlanguage{english} en tête de corps.
\usepackage[english,french]{babel}
\usepackage{fontspec}
\usepackage{geometry}
\geometry{a4paper,margin=25mm,marginparwidth=28mm}
\usepackage{marginnote}
\usepackage{xcolor}
% ulem, not soul. `soul`'s \st and \ul are built on a character-by-character
% scan that XeLaTeX's font machinery breaks: the document fails with an
% undefined control sequence rather than degrading. ulem does the same two jobs
% and is Unicode-engine safe. `normalem` stops it hijacking \emph, which the
% transcriptions use for Grothendieck's own emphasis.
\usepackage[normalem]{ulem}
\usepackage{hyperref}
\hypersetup{colorlinks=true,linkcolor=black,urlcolor=black,pdfborder={0 0 0}}

% --- Métadonnées du lot -----------------------------------------------------
% Reprises telles quelles par le script de rendu, qui les affiche en tête de
% la vue de lecture. Elles fixent ce que le fichier prétend couvrir : sans
% elles, une transcription flotte sans point d'attache dans un fonds de seize
% mille feuillets.

\newcommand{\folder}[1]{\gdef\grothFolder{#1}}
\newcommand{\batch}[1]{\gdef\grothBatch{#1}}
\newcommand{\leaves}[2]{\gdef\grothFirst{#1}\gdef\grothLast{#2}}
\newcommand{\foldertitle}[1]{\gdef\grothFoldertitle{#1}}
\gdef\grothFolder{?}\gdef\grothBatch{?}\gdef\grothFirst{?}\gdef\grothLast{?}\gdef\grothFoldertitle{}

% --- Le résumé --------------------------------------------------------------
%
% Il ouvre la lecture modernisée, sous le titre « Résumé », et il est composé
% en petit. Ce n'est pas un ornement : c'est le seul passage du document qui
% ne soit pas des mathématiques, et un lecteur doit pouvoir le reconnaître
% avant de l'avoir lu — puis le sauter s'il connaît déjà le sujet, ou s'y
% arrêter s'il ne le connaît pas. Le filet qui le suit marque la reprise du
% corps.

\newenvironment{resume}
  {\small\setlength{\parskip}{0.45em}}
  {\par\medskip\hrule\medskip}

% --- Mots-clés ---------------------------------------------------------------
%
% En anglais, en clôture du résumé de la lecture modernisée : le vocabulaire
% moderne sous lequel chercher ce que les feuillets construisent. Le manifeste
% du site (`npm run manifest`) les extrait pour étiqueter la cote entière —
% cette ligne est la source unique des tags, il n'y a pas de fichier à côté.
\newcommand{\keywords}[1]{\par\smallskip\noindent
  {\footnotesize\textsc{Keywords} --- \textit{#1}}\par}

% --- L'appareil critique ----------------------------------------------------

% Le numéro de feuillet, en marge. C'est l'ancre qui permet de régler un
% désaccord sur une formule en regardant *un* feuillet plutôt qu'une chemise
% de six cents. Le site s'en sert aussi pour tourner les pages du fac-similé
% quand on fait défiler la transcription.
\newcommand{\leaf}[1]{%
  \marginnote{\footnotesize\color{gray!70}\textsf{#1}}%
  \label{leaf:#1}\ignorespaces}

% Illisible : on ne devine pas. Un mot inventé qui se lit comme les autres est
% le pire résultat possible.
\newcommand{\ill}{\textcolor{red!60!black}{[\,\ldots\,]}}

% Lecture douteuse : le mot est donné, et signalé comme incertain.
\newcommand{\uncertain}[1]{\uline{#1}}

% Ajout de l'éditeur — une lettre manquante, un mot sous-entendu.
\newcommand{\add}[1]{\textcolor{blue!50!black}{[#1]}}

% Biffé par Grothendieck lui-même. Ce qu'il a rayé fait partie du document :
% une rature dit souvent où la pensée a changé de direction.
\newcommand{\struck}[1]{\sout{#1}}

% Note du transcripteur, sur l'état matériel du feuillet.
\newcommand{\note}[1]{\par\smallskip{\small\color{gray!60!black}\textit{[#1]}}\par\smallskip}

% Note marginale de Grothendieck, distincte du corps du texte.
\newcommand{\marginal}[1]{\par\smallskip\noindent\hspace{1em}%
  {\small\color{gray!60!black}#1}\par\smallskip}

% --- Titre ------------------------------------------------------------------

\AtBeginDocument{%
  \begin{center}
    {\large\bfseries Fonds Alexandre Grothendieck --- cote n\textsuperscript{o}~\grothFolder}\\[2pt]
    {\itshape\grothFoldertitle}\\[4pt]
    {\small feuillets \grothFirst--\grothLast\ (lot \grothBatch)}\\[2pt]
    {\footnotesize\color{gray!60!black}Transcription établie d'après le fac-similé numérisé
     par l'Université de Montpellier.\\
     Les lectures incertaines sont soulignées, les illisibles notées [\,\ldots\,],
     les ajouts entre crochets.}
  \end{center}
  \vspace{1em}}

\endinput


\folder{161-2}
\batch{3}
\pages{41}{60}
\dating{[vers 1963-1973]}
\watermark{Édition de démonstration}
\foldertitle{Algèbre universelle [ou catégories] : notes manuscrites (s.d.)}

\begin{document}

\page{41}
\note{la liste des axiomes sur les anneaux, ouverte p. 38, continue ; la
numérotation est la sienne, et elle est reprise plus bas : le 12) de ce
feuillet devient le 13) de la p. 42}
\[
  X_{\mathrm{cpct}} \simeq \tilde{R}_{\mathrm{top\,cpct}}
\]
cpct est la topologie sur $R$ pour laquelle $X_{i} \to X$ est couv. ssi pour
tt anneau artinien \add{local} \struck{\ill{}} \add{(ou anneau de dim $0$)}
$\mathcal{O}$, $X_{i}(\mathcal{O}) \to X(\mathcal{O})$ est épimorphique.

\marginal{encore un topos non cohérent (\uncertain{bien} que sous-topos d'un
topos cohérent)}

\begin{enumerate}
  \item[11)] \underline{Corps} $=$ \struck{\ill{}} Anneaux réduits (local de
    dim $0$) \note{la parenthèse et « réduit » sont portés au-dessus d'un
    premier membre de droite biffé}
    \[
      A^{\ast} \amalg A \xrightarrow{\ \mathrm{incl},\ 0_{A}\ } A
      \quad \text{est épi}
    \]
    \note{le $\amalg$ est marqué $\hat{R}$ en dessous : la somme est prise
    dans $\hat{R}$}
    \[
      X_{\mathrm{ps.corps}} \simeq \tilde{R}_{\mathrm{top\,\struck{\ill{}}}}
    \]
    corps, top. sur $R$ où les familles couvrantes sont \struck{\ill{}}
    celles qui ont des sections sur chaque fibre.
    \[
      \simeq \tilde{R}_{\mathrm{cons}}
    \]
    $R_{\mathrm{cons}}$ la catégorie des schémas compacts réduits associés aux
    schémas de t.f. sur $\mathbb{Z}$, avec top. de Zariski.
    $\tilde{R}_{\mathrm{cons}}$ formé des foncteurs contravariants qui
    commutent aux \uncertain{lim. proj.} finies.
  \item[12)] corps $=$ compact et \struck{\ill{}} intègre $=$ ps-corps
    intègre \add{(clair)} $=$ ps-corps local \struck{\ill{}} (à prouver)
    \note{deux lignes biffées et repassées entre les deux derniers membres}
\end{enumerate}

\note{le reste du feuillet est barré de traits diagonaux}

\textbf{N.B.}\quad Les ss-topos définis dans 1) à 11) \struck{\ill{}} ont des
pts, donc sont définis par leurs points. Il n'est pas clair qu'il en soit
\uncertain{ainsi} pour les intersections de plusieurs de ces sous-topos. Donc
pour vérifier que deux intersections sont égales, \struck{il \ill{}} ou ont
des rel. d'inclusion déterminées, il n'est pas clair qu'il suffise de le
vérifier pour leurs pts, — i.e. pour les structures ensemblistes. — Sauf si on
vérifie à chaque fois que les int. précéd. ont assez de pts ! Vérifier donc
qu'un anneau \struck{\ill{}} \add{du cpct intègre} \uncertain{dès que} $A$ est
ps corps local, $A$ est \struck{\ill{}} \uncertain{loc. \ill{} intègre}.

\uncertain{2)} La définition \struck{\ill{}} \add{$A : R \to \mathcal{E}$
transforme sommes \struck{\ill{}} finies, et} qu'il transforme familles
couvrantes pour « top. cpct » en familles couv. Montrons qu'une famille
$X_{i} \to X$ qui \struck{\ill{}} est couv. pour \uncertain{top. discrète}
\ill{} de tt (les $X_{i}(k) \to X(k)$ est épi pour tt corps $k$) \struck{est
\ill{}} \ill{} épi. Réc. sur $\dim X$ assez évident.

\page{42}
\begin{enumerate}
  \item[12)] \underline{Anneau compact} (i.e. à spectre compact) —
    $X_{\mathrm{cpct}} \simeq \tilde{R}_{\mathrm{top\,cpct}}$
    \marginal{$\forall U$, $A(U)$ à spectre cpct i.e.}
    \[
      \forall f \in A(U),\ \exists (!)\ e \in A(U),\ e^{2} = e,\
      ef \text{ inversible dans } eA(U),\ \struck{e}\,(1-e)f \text{ nilpotent.}
    \]
    Suffit de le faire pour $f$ universel, $U = A$. On trouve, si
    $D_{n} \subset A \times A \times A$ formé des $g, f, e$ avec $e^{2} = e$,
    $f(1-e) = 0$, $g(1-e) = 0$, $fg = e$, un \uncertain{mono} induit par
    \struck{\ill{}} \uncertain{proj}
    \[
      i_{n} : D_{n} \hookrightarrow A
    \]
    et la condition est que les $D_{n}$ couvrent $A$.

    Pour la top. de $R_{\mathrm{top\,cpct}}$, sera donc la condition que, si
    $D_{n} = \mathrm{Im}\ \amalg_{R} \mathbf{A}^{\ill{}} \xrightarrow{\ i_{n}\ } A$
    (attention, ceci une somme \underline{dans $R$}, pas $\hat{R}$) la famille
    des $i_{n}$ soit couvrante.
    \note{l'exposant de $\mathbf{A}$ se lit $6$ ou $3$ ; $D_{n} \subset A
    \times A \times A$ plaide pour $3$}

    \marginal{topologie non cohérente ! A-t-elle assez de pts ?? — le tout
    encerclé, avec un point d'interrogation en dehors}
  \item[13)] \underline{Anneau compact réduit} $=$ pseudo\struck{corps}
    \marginal{$\forall U$, $A(U)$ \struck{\ill{}} ps. cpct i.e.}
    \[
      \forall f \in A(U),\quad A(U) \xrightarrow{\ \sim\ } A_{f} \times A/fA,
      \text{ i.e. } \exists (!)\ e \in A(U),\ e^{2} = e,\ ef \text{ inv.},\
      \uncertain{e(1-f) = 0}
    \]
    i.e. $i_{1} : D_{1} \hookrightarrow A$ est un isom
    \[
      X_{\mathrm{cpct}} \simeq \tilde{R}_{\mathrm{top\,ps\,corps}}
      \simeq \hat{R}_{\mathrm{cons}}
    \]
    \note{l'indice de $X$ se lit comme au 12) ; le contexte demande « compact
    réduit »}
    ps. corps est la topologie engendrée par les $X' \to X$ qui ont une section
    fibre par fibre ($(X_{i} \to X)$ couvrante ssi $\exists i_{0}$,
    $X_{i_{0}} \to X$ ainsi). $R_{\mathrm{cons}}$ comme avant, localisé de
    $R$ pour les \struck{topologie} flèches \struck{bijections} à ext.
    résiduelles triviales.
    \[
      \hat{R}_{\mathrm{cons}} \subset \hat{R}_{\uncertain{\mathrm{cpct}}}
      \quad \text{formé des } R^{\circ} \to \mathrm{Ens} \text{ qui
      transforment les flèches précédentes \ill{}}
    \]
  \item[14)] \underline{Corps sép. clos} $=$ corps str. hensélien
    \[
      X_{\mathrm{corps\,hens}} \simeq \tilde{R}_{\mathrm{top\,corps\,sépclos}}
    \]
    top. corps hens., topologie où les familles couvrantes sont celles
    \struck{pour \ill{}} $X_{i} \to X$ qui \struck{sont \ill{}}
    \struck{\ill{}} ont \ill{} pt \struck{\ill{}} \uncertain{séparable}
    \ill{} \add{(qui \ill{})} sur chaque \ill{}
    \[
      \simeq \widetilde{(R_{\mathrm{cons}})}
    \]
    et $R_{\mathrm{cons}}$ comme dans 11), avec top. \underline{étale}.
\end{enumerate}

\page{43}
\begin{enumerate}
  \item[15)] \underline{Corps alg. clos} $=$ corps \struck{\ill{}} ultralocal
    \[
      X_{\mathrm{corps\,alg\,clos}} \simeq
      \tilde{R}_{\mathrm{top\,corps\,alg\,clos}}
    \]
    topologie où $X_{i} \to X$ est couvrante si elle est surjective
    \[
      \simeq \tilde{R}_{\mathrm{cons},\,\mathrm{fppf}}
    \]
    $R_{\mathrm{cons}}$ comme dans 11), topologie fppf $=$ topologie des
    familles \struck{contenant une \ill{}} \struck{surjectives} familles
    finies surjectives.
  \item[16)] \underline{Anneau \struck{\ill{}} annulé par $n$}
    (\struck{\ill{}} $\uncertain{\geqslant 1}$)

    (a) $n 1_{A} = 0$, ou $V(n 1_{A}) \xrightarrow{\ \mathrm{isom}\ }
    e_{\varepsilon}$
    \[
      X_{\struck{n\text{-tors}}\ \mathbb{Z}/n\mathbb{Z}}
      \simeq \tilde{R}_{\mathrm{top}\ n\text{-tors}}
      \simeq \hat{R}_{\mathbb{Z}/n\mathbb{Z}}
    \]
    top $n$-tors topologie où $X_{i} \to X$ est couvrante ssi \uncertain{un
    des} $V(n 1_{X_{i}}) \to V(n 1_{X})$ \struck{\ill{}} \add{a une} section
    ($X_{i} \otimes_{\mathbb{Z}} \mathbb{Z}/n\mathbb{Z} \to X
    \otimes_{\mathbb{Z}} \mathbb{Z}/n\mathbb{Z}$) ;
    $R_{\mathbb{Z}/n\mathbb{Z}}$ \struck{top. \ill{}} sous cat. pl. de $R$
    formée des \struck{\ill{}} $X$ qui sont des
    $\mathbb{Z}/n\mathbb{Z}$-schémas ($R_{\mathbb{Z}/n\mathbb{Z}}$ peut
    être considéré comme une catégorie de fractions de $R$ pour les flèches
    $f$ telles que $f \otimes_{\mathbb{Z}} \mathbb{Z}/n\mathbb{Z}$
    bijections).

    \marginal{\struck{\ill{}} sous-topos ouvert de $X_{\mathrm{an}} =
    \hat{R}$ ; N.B. $X_{\uncertain{1}\text{-tors}} = X_{\emptyset} =$ topos
    ponctuel — le tout encerclé}

    \underline{Ex} Si $n$ premier, \ill{} \ill{} les anneaux \ill{}
    \uncertain{corps}.

    (b) anneau où $n$ inversible.
    \[
      X_{\mathbb{Z}[1/n]} \simeq \tilde{R}_{\mathrm{top}\,\mathbb{Z}[1/n]}
      \simeq \hat{R}_{\mathbb{Z}[1/n]}
    \]
    top $\mathbb{Z}[1/n]$, top. où $X_{i} \to X$ couvrantes ssi un des
    $X_{i} \otimes_{\mathbb{Z}} \mathbb{Z}[1/n] \to \struck{X}
    \otimes_{\mathbb{Z}} \mathbb{Z}[1/n]$ a une section ;
    $R_{\mathbb{Z}[1/n]}$, sous-cat. pleine de $R$ des $\mathbb{Z}[1/n]$
    schémas aff. de t.f. \marginal{sous-topos ouvert — encerclé}

    (c) Soit plus généralement $S \to e_{R}$ un mono dans $\mathrm{Pro}\,R$
    ($R = \mathrm{Sch}\text{-}\mathrm{aff}$ ; i.e. $\mathbb{Z} \to \Lambda$
    \uncertain{épi \ill{}}), d'où $R_{/S} \subset R$ sous-catégorie pleine de
    $R$ (correspondant aux anneaux de t.f. sur $\mathbb{Z}$ qui sont des
    $\Lambda$-algèbres). Considérons les Anneaux du type $\varepsilon$ qui
    sont $\Lambda$-algèbres (la structure de $\Lambda$-algèbre est-elle
    unique ? Oui, on est ramené par le procédé standard au cas ensembliste) ;
    les $\Lambda$-algèbres de $\varepsilon$ sont les foncteurs \ill{}
    $\varepsilon \to \struck{\mathrm{Ens}}\ (\Lambda\text{-alg})
    \longleftrightarrow (\mathrm{Ann})$ qui sont représentables
    ensemblistement. Or \struck{\ill{}} \struck{$(\Lambda\text{-alg}) \simeq
    \underline{\mathrm{Hom}}$} Or s'identifient (\struck{Ann}) \add{Ann à}
    $\underline{\mathrm{Hom}}'(R, \mathrm{Ens})$
    \note{de (c) à la fin du feuillet, le texte est barré de traits
    diagonaux ; il continue sans rupture p. 44, qui n'est pas barrée}
\end{enumerate}

\page{44}
(où $'$ signifie exact à g.) et $\Lambda$-alg. à
$\underline{\mathrm{Hom}}'(R_{\Lambda}, \mathrm{Ens})$ — $R_{\Lambda}$ cat.
des schémas de t.f. sur $\Lambda$, le foncteur d'inclusion
$\underline{\mathrm{Hom}}'(R_{\Lambda}, \mathrm{Ens}) \hookrightarrow
\underline{\mathrm{Hom}}'(R, \mathrm{Ens})$ est transposé des

\begin{tikzcd}
\underline{\mathrm{Hom}}'(R_{\Lambda}, \mathrm{Ens}) \arrow[r, hook] &
\underline{\mathrm{Hom}}'(R, \mathrm{Ens}) \\
(\Lambda\text{-alg}) \arrow[u, "\wr"] \arrow[r, hook] &
(\mathrm{Ann}) \arrow[u, "\wr"']
\end{tikzcd}

avec, à gauche, $A \mapsto \bigl(B \mapsto \mathrm{Hom}(B, A)\bigr)$
($B$ une $\Lambda$-alg), et à droite $A \mapsto \bigl(B \mapsto
\mathrm{Hom}_{\mathbb{Z}}(B, A)\bigr)$ ($B$ une $\mathbb{Z}$-alg
\struck{$\Lambda$-alg}), $\simeq \mathrm{Hom}_{\Lambda}(B_{\Lambda}, A)$ si
$A$ une $\Lambda$-alg.

foncteur $R \to R_{\Lambda}$, $X \mapsto X \otimes_{\mathbb{Z}} \Lambda$
i.e. $X \mapsto X \times_{e} S = X \times S$ (qui est en effet exact à g.).
Donc
\[
  X_{\Lambda\text{-alg}} \simeq \tilde{R}_{\Lambda\text{-top}}
  \simeq \hat{R}_{\Lambda}
\]
$\Lambda$-top : $X_{i} \to X$ couvrantes ssi les $X_{i} \struck{\times} \times_{e}
S \to X \times_{e} S$ ont une section. \textbf{N.B.}\ $R_{\Lambda}$ est aussi
déduit de $R$ par qu. cat. de fractions, en rendant les $f$ telles que
$f \times_{e} S$ un isom.

\underline{Ex}\quad $X_{\mathbb{Z}/n\mathbb{Z}\text{-alg}}$,
$X_{\mathrm{car}\,p}$ (cas particulier), $X_{\mathbb{Z}[1/n]}$ ;
\[
  X_{\mathbb{Q}} = X_{\mathrm{car}\,0} = \struck{\ill{}}
  \bigcap_{n \geqslant 2} X_{\mathbb{Z}[1/n]}
\]

\textbf{N.B.}\quad Ce qui vient de \add{se} voir se généralise au cas d'une
\ill{} finies, et l'\ill{} de structure donnée $\Lambda \in T(\mathrm{Ens})
= \mathrm{Ind}\,R^{\circ}$, \ill{} suppose que $S \to e_{R}$ \ill{}, où
$R_{S} =$ catégorie des foncteurs \ill{} que déduit des isom pour \ill{} image
ess. de $\mathrm{Pro}\,R_{/S}$, $X \mapsto X \times_{e} S$ \struck{\ill{}}.
Comme $S \in \mathrm{Ob}\,\hat{R}$ (\struck{$S \subset e_{\hat{R}}$}) \ill{}
$\mathrm{Pro}\,R \hookrightarrow \hat{R}$, \struck{\ill{}} \ill{} dans
l'image essentielle de $\mathrm{Pro}\,R$, et s'intéresse aux $f :
\varepsilon \to \hat{R}$ tels que $f^{*}(S) \to f^{*}(e_{\hat{R}}) =
e_{\varepsilon}$ \ill{} \ill{} $f^{*}(S) \hookrightarrow e$ \ill{} \ill{}
des \ill{}-structures qui sont \uncertain{représentées} \ill{} par des
\underline{ouverts} du topos classifiant $\hat{R}$.

\marginal{$R$ quelconque avec $\varprojlim$ finies \uncertain{commutant}
\ill{}. On se donne i.e. $S \in \mathrm{Pro}\,R$, un mono, et on considère
\ill{} de $R$ pour les flèches que \ill{} $\times_{e} S$, donc sous-catégorie
\ill{} pl. $R \to \mathrm{Pro}\,R_{/S}$, $X \mapsto$ \ill{}}

\marginal{$\pi(\mathrm{Ens})$, \ill{} — montrons que si $S \in
\mathrm{Pro}\,R$ \ill{} les \ill{} conditions \ill{}}

\page{45}
\note{en tête de feuillet, à droite : « 16 (suite) »}
($\forall$ cat. $R$ stable par $\varprojlim$ finies \add{ess. petite}) Plus
gén., pour tout mono $S \xrightarrow{\ i\ } e_{R}$ dans $\hat{R}$, d'où
sous-catégorie $R_{/S}$ de $R$, \struck{Est sous-catégorie} considérons
(\struck{les} topos $\varepsilon$-\ill{}) les $R \xrightarrow{\ \varphi\ }
\varepsilon$ ex. à g. \struck{tels que} \add{qui \ill{} $\varphi(i)$ un isom.}
i.e. $\forall X \in \mathrm{Ob}\,\varepsilon$, $\mathrm{Hom}(X, -) \circ
\varphi : R \to (\mathrm{Ens})$ ex. à g. \add{(avec \uncertain{versions}
\ill{} \ill{})}. Or les foncteurs ex. à g. \struck{\ill{}} $R \to
\mathrm{Ens}$ s'identifient \ill{} aux éléments de $\mathrm{Pro}\,R$
(\ill{} pro-représentables), et \struck{\ill{}} $\xi \in \mathrm{Pro}\,R$
tels que $\mathrm{Hom}_{\mathrm{Pro}R}(\xi, S) \neq \emptyset$, on prouve dans
les \ill{} $S$, donc en sont ass. les \ill{} i.e. \ill{} ou encore la
sous-cat. pleine de $\mathrm{Pro}(R_{/S})$, \struck{donc} de
\[
  \mathcal{C} = (\mathrm{Pro}(R))^{\circ} = \mathrm{Ind}(\mathcal{S})
  \qquad (\struck{\ill{}} = R^{\circ})
\]
(qui sont \struck{\ill{}} \ill{}) tel qu'il existe un morphisme $\Lambda \to
\xi$, i.e. $\Lambda \to$ \ill{} de $\mathcal{S}$ correspondant à $S$. On a
pour les topos \struck{\ill{}} classifiants \ill{} : $\hat{R}_{/S}$, avec
$\hat{R}_{/S} \xrightarrow{\ \hat{\imath}\ } \hat{R}$ déduit de $i :
R_{/S} \to R$. On a : $i^{*} =$ inclusion de $R_{/S}$ en $R$,
$i^{*}(\struck{X}) = X \times_{e} S / S$, $i_{*}(Y)$ \add{(objet de
$R_{/S}$)} $= \{\,X \mapsto \mathrm{Hom}_{S}(X \times S, Y) =
\mathrm{Hom}(X \times S, Y)\,\}$, $X \in \mathrm{Ob}\,R$, i.e.
$i_{*}(\struck{Y}) = \underline{\mathrm{Hom}}(S, Y)$.

Considérons maintenant une famille $S_{\alpha} \xrightarrow{\ i_{\alpha}\ }
e_{B}$ \add{filtrante décroissante} de ss-objets de $e_{R}$, et soit $T' =
\bigcap_{\alpha} T_{S_{\alpha}}$ l'\ill{} ss-espèce de structure de $T =
T_{R} = T_{\hat{R}}$ définie par les conditions $\varphi(i_{\alpha})$
\add{$= f^{*}(i_{\alpha})$}, i.e. $f : \varepsilon \to B = \hat{R}$)
\uncertain{iso} pour tout $\alpha$. Définie aussi par les conditions sur les
$\mathrm{Hom}(X_{\xi}, -) \circ \varphi$, on \add{est} ramené à examiner le
cas du topos ponctuel $\varepsilon = \mathrm{Ens}$ : \struck{Oui \ill{} donc
\ill{}} Soit $S = (S_{\alpha}) \in \mathrm{Pro}\,R$, \struck{$S \to e_{R} =
e_{\mathrm{Pro}R}$ un mono} d'où $\xi = (\xi_{\alpha}) \in
\mathrm{Ind}\,\mathcal{S}$, les objets \ill{} $\mathcal{O}$ de $\mathcal{S}$
\ill{} sont ceux pour lesquels il existe un (unique) morphisme $\xi \to
\mathcal{O}$, donc dans $\mathrm{Pro}\,R = \mathcal{R}$ un \add{(unique)}
$x \to S \hookrightarrow e_{R} = e_{\mathcal{R}}$ \ill{} \ill{}, qui \ill{}
\ill{} pour un \ill{}. Mais quelles sont les catégories classifiantes ?
\struck{C'est le \ill{} topos} \add{Le topos $\hat{R}'$ est le \ill{}} \ill{}
ces \ill{}, mais quels sont ces pts \ill{}. \ill{} $\mathcal{R}'$ \add{une
petite \ill{}} avec $\varprojlim$ \ill{} \ill{} $\bigcap_{\alpha}
\hat{R}_{/S_{\alpha}}$. On a déterminé
\note{le feuillet s'arrête ici en cours de phrase ; la marge, écrite le
cahier tourné, porte le développement qui suit}

\marginal{$\to \mathrm{Pro}\,R_{/S}$ définis par objets de base. On trouve
donc \add{petite} la catégorie \uncertain{modulaire} avec les conditions et
$\varprojlim R_{/S_{\alpha}}$ qui est une catégorie de fractions de $R$ pour
les flèches $i$ telles que $i \times_{e} S$ un iso. \ill{} des
\uncertain{ind-représentables} des $R \to \mathrm{Pro}$ \ill{} limites de la
catégorie $\varprojlim_{\alpha} R_{/S_{\alpha}}$.
\[
  \hat{R}_{/S_{\alpha}} = \bigcap_{\alpha} \hat{R}_{/S_{\alpha}}
  = \widetilde{\varprojlim_{\alpha} R_{/S_{\alpha}}}
\]
plongés dans $B$ par $i_{\alpha *}$, un $i_{\alpha !}$ !! \ill{} sauf erreur,
ce sont les \ill{} $X \to X \times_{e} S$, on \ill{} les \ill{}
$\bigcap_{\alpha} B_{\bar{I}_{\alpha}} = \bigcap_{\alpha}$ \ill{}}

\page{46}
\begin{enumerate}
  \item[17)] \underline{Anneau parfait de car. $p$} ($p > 0$)
    \[
      p\,1_{A} = 0, \qquad A \xrightarrow{\ x \mapsto x^{p}\ } A \text{ isom}
    \]
    \marginal{en car. $0$, \uncertain{pas} \ill{} \uncertain{réduit} ? \ill{}
    \ill{} — écrit en biais}

    \textbf{N.B.}\ parfait $\Longrightarrow$ réduit
    \[
      X_{p\text{-parf}} = \tilde{R}_{\mathrm{top}\,\struck{\ill{}}}
      = \hat{R}_{p\text{-parf}}
    \]
    $R_{p\,\mathrm{parf}}$ catégorie des schémas affines \add{parfaits} de
    car. $p$ qui sont clôtures parfaites des schémas \uncertain{de t.f.} sur
    $\mathbb{Z}/p\mathbb{Z}$.

    top. $p$-parf sur $R$ : $X_{i} \to X$ est couvrante si $\exists i_{0}$,
    $(X_{i_{0}})_{p} \to (X)_{p}$ admet une section après passage à la
    clôture parfaite de $X_{p}$.
  \item[18)] \struck{Anneau \ill{} corps} \add{compact} parfait de car. $p$
    ($> 0$). parfait de car. $p$ $\cap$ ps. corps.

    top. cpct $p$-parfaite : $X_{i} \to X$ est couvrante ssi $\exists i_{0}$,
    $X_{i_{0}} \to X$ a un pt \uncertain{radiciel} au dessus de \ill{} pt de
    $X_{p}$
    \[
      X_{\struck{\ill{}}\ p\text{-parf}} \simeq
      \tilde{R}_{\struck{\ill{}}} \simeq \hat{R}_{\mathrm{cons}\ p\text{-parf}}
    \]
    $R_{\mathrm{cons}\ p\text{-parf}} =$ catégorie des clôtures parfaites des
    schémas de car. $p$ dans $R_{\mathrm{cons}}$.
  \item[19)] \underline{Corps parfait} \struck{\ill{}} \add{(car $p > 0$)}
    \[
      X_{\mathrm{corps}\ p\text{-parf}} \simeq
      \tilde{R}_{\mathrm{top\ corps}\ p\text{-parf}} \simeq
      \tilde{R}_{\mathrm{cons}\ \uncertain{p\text{-parf}}\ \mathrm{zar}}
    \]
    $X_{i} \to X$ est couvrante ssi $X_{i,p} \to X_{p}$ \add{a} des pts
    radiciels \ill{} \uncertain{sur chaque} fibre.
  \item[20)] \struck{Anneau à caractéristique} \struck{Anneau parfait (de
    car. \ill{} parfaite, \ill{} de car. $p > 0$ un précisé (inclus par
    $0$)}
    \note{les deux lignes du titre sont biffées l'une après l'autre ; le
    reste du feuillet est barré de traits diagonaux}

    \struck{A.} On veut que loc sur $\varepsilon$ il ait les propriétés
    envisagées. Soit $C$ l'un des \ill{} qui \ill{} \ill{} \ill{}
    (\uncertain{comme} cas de car. $0$), \struck{Donc} soit \ill{} l'ensemble
    \ill{}. Alors $\varepsilon$ \ill{} \ill{} \ill{} \ill{} \ill{}
    \[
      e_{\varepsilon} \simeq \coprod_{c \in C} U_{c},
    \]
    un $U_{c}$ de car. $p$ si $c \neq \emptyset$, et $\Delta U_{\emptyset} = 0$.
    Puis alors $X_{\uncertain{\mathrm{parf}}} \subset X_{\mathrm{ann}}$ ; leur
    intersection \ill{} \ill{}
    \note{les deux dernières lignes sont surchargées et se lisent mal}
\end{enumerate}

\page{47}
\begin{enumerate}
  \item[20)] \underline{Anneau à car.}

    loc. sur $\varepsilon$, $A$ est de car $p > 0$ ou de car. $0$, i.e. la
    famille \struck{\ill{}}
    \[
      \Bigl\{\, V(p \cdot 1_{A}) \hookrightarrow e_{\varepsilon} \ ;\
      \bigcap_{n \geqslant 1} (e_{\varepsilon})_{n 1_{A}} \hookrightarrow
      e_{\varepsilon} \,\Bigr\}
    \]
    est épimorphique (cela fait \textbf{factoriser} \add{univers} pour le
    sous-topos $\Sigma_{\mathrm{car}}$ \ill{} $\varepsilon \to \hat{R}$ \ill{}
    une $\varprojlim$ \underline{infinie} \ill{}, donc ne donne pas
    l'existence d'un topos classifiant). On aura $\varepsilon \to \hat{R}$ se
    factorise localement sur $\varepsilon$ par la « réunion » $V$
    \add{$= X_{\mathrm{car} > 0}$} des sous-topos ouverts
    $\hat{R}_{\mathbb{Z}/p\mathbb{Z}} \simeq
    \hat{R}_{\mathrm{Sp}\,\mathbb{Z}/p\mathbb{Z}}$, et le sous-topos $W =
    \hat{R}_{\mathbb{Q}}$ (qui n'est \uncertain{pas} sous-topos ouvert) avec
    \ill{} \underline{images}.
    \[
      X_{\mathrm{car}} = \tilde{R}_{\mathrm{top.car}}
    \]
    \struck{\ill{}} Où « top car » est la borne inférieure des « top car $p$ »
    pour \ill{} \ill{} premiers $p$ et pour \ill{}, i.e. les familles $X_{i}
    \to X$ \ill{} couvrantes ssi pour \struck{\ill{}} tt $p$, il existe un
    indice $i(p)$ tel que $(X_{i(p)})_{p} \to (X)_{p}$ (fibres au dessus de
    $p$) ait une section. C'est [p. ex. les familles des $(X_{p} \to X)_{p}$
    \ill{} \ill{} \add{finies} $\to$ couvrantes, et si $X_{\mathbb{Q}} \neq
    \emptyset$, il n'y a pas de familles couvrantes finies
    \uncertain{cofinales}]

    \marginal{\underline{Topos non cohérent} — encerclé, avec un croquis :
    deux ovales $Z'$, $Z''$ se recouvrant, la région commune marquée $I$, une
    flèche vers le bas}

    Plus généralement, soit $B$ un topos, $V$ et $W$ deux sous-\ill{},
    considérons $T = T_{B}$, et $T_{V}$ et $T_{W}$, \ill{} à foncteurs
    \underline{localement} par $B$ \ill{} $T' = (\varepsilon \to
    \struck{f} \varepsilon \to B$, \ill{} $B \simeq B'$. \ill{} \ill{} \ill{}
    $f$ que \struck{l'image} $e_{\varepsilon} = \sup(e', e'')$, i.e. c'est les
    plus \ill{} \ill{} de $\varepsilon$ \struck{\ill{}} \add{tel pt} \ill{}
    comme $f|_{\varepsilon_{/e'}}$ se factorise par \ill{}, et pour $e''$. On
    a alors $f|_{e' \wedge e''} \to V \wedge W$, \ill{} \ill{} plus quel
    ouvert \struck{\ill{}} $\subset e_{\varepsilon}$ \ill{} \struck{\ill{}}
    est aussi \ill{}. \ill{} $e' \wedge e''$ \ill{} \ill{} ? [Si \ill{}
    $U_{c}\,e'$, $U_{c}\,e''$ dans $U_{c}\,e' \wedge e''$, OK] / Soit \ill{}
    un topos, un sous-topos \struck{\ill{}} $Z \subset B$ \struck{sous-topos}
    contenant $V_{\varepsilon}$ et $W$ et ouvert $Z' = Z_{/e'}$, $Z'' =
    Z_{/e''}$, et \struck{\ill{}} $\sup(e', e'') = e_{Z}$ i.e. \ill{} des
    sous-topos de $Z$, ou de $B$, \ill{} \ill{} pour $T$, c'est \ill{} \ill{}
    tels que $Z$ \struck{\ill{} sous-topos} ayant des sous-topos ouverts avec
    $Z' \subset V$, $Z'' \subset W$, et $Z = \sup(Z', Z'')$ (dans l'un de
    \ill{}, c'est kif kif). Ou \ill{} $Z = \sup(Z', Z'') \subset \sup(V, W)
    \subset Z$, donc $Z = \sup(V, W)$. De plus on aura \struck{\ill{} $Z'$
    \ill{} $Z''$} $Z' = Z \cap V' = Z \times_{Z} V =$ ouvert de $V$, $Z'' =$
    ouvert de $W$.
    \note{les deux colonnes du bas de page s'imbriquent ; l'ordre ci-dessus
    suit les renvois}

    \marginal{note en biais dans la marge, largement illisible : \ill{} $Z$
    \ill{}}
\end{enumerate}

\page{48}
\note{en tête de feuillet, à droite : « (20 suit) »}
\struck{\ill{}} $V \wedge W = U$ ; d'ailleurs \add{$Z'$ et $Z''$ suivant
$I$} (ouvert \ill{} $Z'' \wedge V$ un ouvert de $V$, (mais pas \ill{} ouvert
$Z'$, ce ouvert $V$, et \ill{}. Pb ait une solution, il f et s. \struck{$Z'$
de $V$} \struck{\ill{}} ouvert de $V$ qui induise sur $W$, et $Z''$
\struck{\ill{}} induise un ouvert $Z'' \wedge V$ sur $V$, \add{$V = \sup(Z',
Z'' \wedge V)$}, $W = \sup(Z'', Z' \wedge W)$. \ill{} nécessaire. Pour voir
que c'est \ill{} \ill{} le topos $Z^{*}$ qui est obtenu \ill{} suivant leur
intersection $Z_{1}$, $Z \to B$ qui est un plongement \ill{} morphismes $V
\to Z^{*}$ et $W \to Z$ \ill{} $Z' \xrightarrow{\ \mathrm{Id}\ } Z'$ et $V
\wedge Z'' \to Z''$ d'une part \struck{\ill{}} \ill{} d'autre part, \ill{} $Z
\to B$ ou \ill{} (ces \ill{} avec \ill{}), donc $V, W \subset Z \subset B$.
\ill{} que $Z$ satisfait : la condition \ill{} que $Z' \wedge Z''$ soit ouvert
\ill{} $Z''$ \ldots

Trouver donc $Z' \wedge Z''$ est un ouvert (dans $Z$ est obtenu en recollant
\ill{}), De $Z'' \wedge V$ un \ill{} \add{par des $Z''$} qui, avec \ill{}
\ill{} en sens ouvert. \underline{Le \ill{}} : Pour que le Pb \struck{\ill{}
un ouvert $Z'$ de} \struck{\ill{} $\exists Z'$ \ill{} le plus} \ill{} ouvert
\struck{\ill{}} $Z' \wedge W$ sur $W$ \ill{} ouvert de $W$ qui induise, alors
$V = \sup(Z', Z'' \wedge V)$. On a vu que c'est nécessaire — suffisant,
considérons \ill{} en recollant $Z'$ et $Z''$ suivant \ill{}, d'où un
morphisme $Z \to$ \ill{} (à vérifier) et des morphismes qui sont définis via
$Z' \to$ \ill{} ($V = \sup(Z', V \wedge Z'')$) et via \ill{} ce sont des
plongements (ces \ill{} \ill{} les plongements donc \ill{}). On voit alors
aussitôt que \ill{}.

\textbf{N.B.}\ Il suffit un pt \ill{} dans $Z'$ et dans $Z''$.

\underline{Dans} le cas envisagé $W = X_{\mathrm{car}\,0}$, \struck{\ill{}}
\ill{} de $B = \hat{R}$ (correspondant \ill{} $X \in \mathrm{Ob}\,R$ qui sont
une \ill{} précision) ; \textbf{N.B.}\ ce crible \ill{} \ill{} des doubles
\uncertain{fibrés} et $V \wedge W$ est l'ouvert \ill{} à
\underline{l'anneau nul} \ill{}. Donc \underline{OK}.

\marginal{en regard, dans la colonne de droite :
\[
  B = \hat{R} = X_{\mathrm{ann}}, \qquad V = X_{(\mathrm{car} > 0)}
\]
$V$ est \uncertain{déjà} un ouvert \ill{} des cribles des objets \ill{} en
positive \struck{\ill{}} \ill{}, ne \struck{\ill{} \ill{} \ill{} produits
finis} \add{produits finis} \ill{} (mais pas d'objet final) ; \ill{}
$X_{\mathrm{car}\,0} \simeq \hat{R}_{\mathbb{Q}}$ correspondant (initial)
\ill{} objet \ill{} de $R_{\mathbb{Q}}$. \textbf{N.B.}\ $X_{\mathrm{car} >
0} \simeq \tilde{R}_{(\mathrm{car} > 0)}$, $R_{\mathrm{car} > 0}$ sous-cat.
pl. de $R$ des schémas \ill{} de car $> 0$ $=$ somme \uncertain{amalgamée}
des $R_{\mathbb{Z}/p\mathbb{Z}}$ ($p$ fini) sous \ill{} $\hat{R} \to
\hat{R}_{\mathbb{Z}/p\mathbb{Z}}$}

\marginal{note en biais, largement illisible : constructions précédentes
\ill{}}

\page{49}
\[
  X_{\mathrm{car}} \simeq \hat{R}_{\mathrm{car}},
\]
où $R_{\mathrm{car}}$ est la catégorie somme amalgamée des $R_{p}$ ($p$
premier \struck{\ill{}}) sous $\struck{\ill{}} \to R_{p}$ \ill{} tt $p$ ;
\add{les des} schémas \struck{\ill{}} qui sont de type fini sur un corps
premier (y inclus $\mathbb{Q}$) \struck{\ill{}}

\begin{enumerate}
  \item[21)] \underline{Anneaux à car. avec propriétés précisées}
    \struck{\ill{}} Re \underline{N.B.}\ On peut reprendre maintenant toutes
    les notions envisagées dans 1) à 19 au cas donnée en propriété plus gén.,
    et \ill{} un sous-topos \struck{$Z(p)$} de $\hat{R}$ \struck{\ill{}
    $\hat{R}_{p}$}, en trouvons comme topos modérément \ill{} sous-topos
    correspondant \ill{} comme ceux des $Z(p)$ tels que \ill{} de \ill{} i.e.
    l'anneau nul \ill{}, si \ill{} $\emptyset_{p} \notin Z(p)$ i.e. pour
    $Z(p)$, et des 2 sommes \add{$\struck{\text{catégorie}}$ topos}
    \uncertain{fibrantes} pour deux topos $Z(p)$ [munis des \ill{}
    amalgamées sous \ill{} ouvert \ill{} $\emptyset_{p}$]. Il des \ill{}
    sous-topos ouverts \ill{} chacun des $Z(p)$, y a assez de pts ssi \ill{}
    a assez \ill{}.
    \note{$Z(p)$ transcrit le signe qui désigne ici le sous-topos de la
    caractéristique $p$}
  \item[22)] \underline{Anneaux quasi\struck{\ill{}} connexes} \add{de rang
    $\leqslant n$} (dans le spectre a au plus $n$ comp. connexes, donc qui
    \ill{} sommes d'au plus $n$ schémas connexes) \struck{A} \struck{Déf} On
    peut trouver loc. \ill{} loc. une décomposition en produit $A_{1} \times
    A_{2} \times \cdots \times A_{n}$, avec les $A_{i}$ connexes
    $\Longleftrightarrow$ on peut trouver loc. \struck{\ill{} des} \ill{}
    projecteurs $e_{i}$ ($1 \leqslant i \leqslant n$) orthogonaux de somme
    $1$, avec $A_{i} = A / \sum_{j \neq i} e_{j} A$ connexe pour tt $i$,
    $1 \leqslant i \leqslant n$.
    \[
      \Pi_{n} \subset A^{n} \ \hookleftarrow\ e_{1}, \ldots, e_{n} \text{ avec }
      \begin{cases}
        e_{i}^{2} = e_{i} \quad \forall i \\
        e_{i} e_{j} = 0 \quad \forall i, j,\ i \neq j \\
        \sum e_{i} = 1
      \end{cases}
    \]

    \begin{tikzcd}[column sep=small, row sep=small]
    A \times \Pi_{n} \arrow[d] & A_{j} \arrow[dl] \\
    A_{\Pi_{n}} \arrow[d, Rightarrow] & \\
    \Pi_{n} &
    \end{tikzcd}

    \[
      A_{j} = A_{\Pi_{n}} \Big/ \sum_{j \neq i} e_{j} A_{\Pi_{n}}
      \qquad
      \Pi_{n} \amalg \Pi_{n} \xrightarrow{\ i_{n,j}\ } \Pi(A_{j})
      \quad \text{(proj. de } A_{j}\text{)}
    \]
    \note{le diagramme est un croquis : $A \times \Pi_{n}$ au-dessus de
    $A_{\Pi_{n}}$, une double flèche vers $\Pi_{n}$, $A_{j}$ tombant sur
    $A_{\Pi_{n}}$ ; les deux flèches obliques de $\Pi_{n} \amalg \Pi_{n}$
    vers $\Pi_{n}$ ne sont pas rendues}
    soit $\Pi_{n,j} \subset \Pi_{n}$ plus grand ss-objet pour lequel
    $i_{n,j}$ un épi. Comment avec l'objet de base \underline{dans
    $\varepsilon$}
\end{enumerate}

\page{50}
\note{en tête de feuillet, à droite : « (22 suit) »}
Soit $\Pi'_{n} = \bigcap_{j} \Pi_{n,j}$, et $e'_{n} = \mathrm{Im}(\Pi'_{n}
\to e_{\varepsilon})$. La condition est que $e'_{n} \hookrightarrow
e_{\varepsilon}$ est un \underline{iso}, i.e. $\Pi'_{n} \to e_{\varepsilon}$
un épi. Il n'est pas clair cependant que la formation des $e'_{n}$ commute
aux changements de pt de topos de base $\varepsilon' \to \varepsilon$, car ce
n'est pas clair pour la formation de $\Pi_{n,j}$. Donc il n'est pas clair
qu'on a un topos \uncertain{modérateur} \struck{\ill{}}, d'abord de pts
\ill{} gens des \ill{} de \underline{plus} de rapport tout. Cependant le
choix des \ill{}, \struck{à cette fois} à permutation près, \struck{\ill{}}
on désigne le $n$-uple des $e_{i}$ qui sont $\neq 0$ \struck{\ill{}} les
$e_{i}$ \ill{} tous différents \ill{} $e_{i}^{2} = e_{i}$, $e_{i} e_{j}
(= e_{i}^{\uncertain{?}}) = 0$ si \ill{} \ill{} \ill{}. Alors on voit que le
choix des $e_{i}$ est bien déterminé, et de plus que si localement \ill{}
\ill{} \ill{} alors \struck{la décomposition} c'est pour ceux $e_{i} =
\ill{}$ typiques $e_{i} = 0$ !) si $n = m$, \ill{} \ill{} que \struck{les
choix des} choix possibles des $e_{i}$ \ill{} est un \ill{}
\add{symétrique} \ill{} $n$ ; donc que \ill{} \ill{} le \ill{}
\underline{\uncertain{lexact}} \ill{} (\ill{} à support \struck{\ill{}}
tout). \struck{\ill{}} $e_{i}$ faisceaux des choix \add{de}
pseudo-torseur \add{sous \ill{}} sous le \struck{\ill{}} $A$ \add{est}
\underline{connexe} \add{\ill{}} typiques. \ill{} définition, \ill{}. Je ne
sais si les \ill{} (mais le torseur \ill{}, \ill{} \ill{} $A$
quasi-connexes de rang exactement $n$ \ill{} tel revient : la donnée
\struck{\ill{}} d'un $\varepsilon \to B$, mais \struck{\ill{}} avec
$\varepsilon$ muni de $A$ \ill{} \ill{} (ou plus des $\varepsilon \to
\hat{R} = B$) ;
\begin{itemize}
  \item[a)] d'un $\mathfrak{S}_{n}$-torseur $P$ (sur $\varepsilon$, i.e.
    d'un morphisme $\varepsilon \xrightarrow{\ P\ } B_{\mathfrak{S}_{n}}$) ;
  \item[b)] d'une \struck{\ill{}}composition des $A_{P}$ en produit
    $\prod_{1 \leqslant i \leqslant n} A_{i}$ avec les $A_{i}$ connexes
    (\ill{} à support $\ill{}$), et avec comp. distincts \ill{} \ill{} et
    l'action de $\mathfrak{S}_{n}$.
\end{itemize}
Cela signifie qu'on a \ill{} le torseur $P$ (i.e. $p$), des anneaux
\underline{connexes} \ill{} à \ill{} \ill{} pour tt $A_{i}$ \ill{} sur $P$,
et \ill{} action de $\mathfrak{S}_{n}$ sur les $P$ permutant \ill{} avec les
$A_{i}$ \ill{} suivant des actions \ill{} sur les indices.

\marginal{le \ill{} à vérifier}

\page{51}
\note{feuillet volant, papier ligné ; la page est une feuille de questions
sur la « typicité » des propriétés de modules et d'algèbres, en plusieurs
blocs disposés librement}

$M$ \underline{plat} cela est-il une propriété typique [elle est stable par
$\varinjlim$ filtr. et par changt de topos]
\[
  N \mapsto M \otimes N \text{ exact}
\]
$\forall f_{1} \cdots f_{n}$ sections de $A$ sur $U$,
\[
  \Bigl(\sum f_{i} A\Bigr) \otimes M \longrightarrow M \quad \text{injectif}
\]
Il suffit de l'exiger pour les \add{$f_{i}$ sect. de $\mathcal{U}$} sections
universelles, i.e. $U = A^{n}$
\[
  \Bigl(\sum t_{i} A_{U_{n}}\Bigr) \otimes M_{U_{n}} \longrightarrow
  M_{U_{n}} \quad \text{injectif.}
\]
$U_{n} = A^{n}$
\[
  Y \longleftarrow X
\]
\[
  M \to M_{i}, \qquad Y \leftarrow Y_{i}, \qquad Z \to Y, \quad N \text{ plat}
  \qquad\qquad M \to M_{i}, \quad A \to A_{i}, \quad N \text{ plat}, \quad B
\]
\note{deux croquis de changement de base, sans flèches complètes ; le second
porte « $N$ plat » sur une verticale au-dessus de $B$}

\marginal{$M$ \underline{fid plat} est-ce typique ? — encerclé}

Donc $\forall N \neq 0$, $M \otimes N \neq 0$

$E \wedge$ \quad $F$ \quad $E_{A}$

$\underline{\mathrm{Modpf}}(A)$ champ des modules \ill{} sur $A$ ;
$\underline{\mathrm{Hom}}\bigl(\underline{\mathrm{Modpf}}(A),
\underline{\mathrm{Fl}}(F)\bigr)$ Catégorie des hom. de champs des \ill{} ;
\struck{Les objets} ; $E \to F$ hom. de champs \underline{constants} ;
$M \mapsto X$ ; d'où $A_{F}$ ; $R$.

\underline{Algèbres} [de t.f. / p.f. / lisses / étales / nettes] pas typiques
\ill{} \ill{} aux \ill{}

\begin{itemize}
  \item \underline{radicielles} / monomorphiques \struck{\ill{}} OK
    (prod. \ill{}) ;
  \item \underline{surjectives} ?! ;
  \item plates \add{OK} / fid plates ;
  \item formellement lisse / — nette OK / — étale ;
  \item entières OK ;
  \item hensé. universel.
\end{itemize}

Considérons une Algèbre qui est déjà de \underline{p.f.} On veut « le
vérifier » (à laquelle lisse etc.). Cela se fait-il pour un sous-topos ??

\page{52}
\struck{\underline{Propriétés relatives d'anneaux}}

propriétés \underline{non typiques} (pour topos) (absolues pour anneaux)
\begin{itemize}
  \item artinien, noethérien, de t.f. sur $\mathbb{Z}$ (ou sur $\mathbb{Q}$,
    \struck{\ill{}} ou $\mathbb{Z}_{S^{-1}}$ \ldots) ;
  \item $n$ composantes connexes ;
  \item $n$ composantes irréductibles.
\end{itemize}

Anneaux plongeables dans des corps : il \underline{faut} qu'ils soient
intègres, mais ce n'est pas suffisant \add{\ill{} s'il \ill{} locale
hypothèse}. Pour l'anneau intègre \underline{universel}, \add{(voir anneau
\ill{} intègre \ill{} \ill{})} les sections locales qui sont partout non
nulles (donc qui sont $A$-régulières) sont celles qui sont inversibles :
l'anneau total des fractions est égal à $A$ lui-même. Une condition
nécessaire pour qu'\struck{\ill{}} un anneau se plonge dans un corps, c'est
que le \add{$X' \to X$} \struck{\ill{}} anneau se plonge dans les sections
nulles (donc \ill{} au \ill{} \ill{} direct de $A$), et c'est sans doute
suffisant

\underline{Modules sur anneaux}
\begin{itemize}
  \item localement libre ; libre, projectif / injectif (?) ;
  \item localement libre de rang $n$ ; libre de rang $n$ ($n \geqslant 1$) ;
  \item de t.f. ;
  \item de prés. finie ;
  \item monogène ;
  \item \struck{\ill{}}
\end{itemize}
\marginal{en regard des trois derniers : ne \uncertain{passent pas à la
lim. ind. filtrante}, donc pas typiques}

\marginal{à gauche, un croquis : $M \to N$ au-dessus de $A \to B$, $f$ et
$g$, $Y \leftarrow X$}

\[
  X_{\mathrm{mod}} = \widehat{\mathrm{Mod}^{\circ}_{\mathrm{pf}}}
\]
$\mathrm{Mod}_{\mathrm{pf}} =$ catégorie des modules de prés. finie sur
$\mathbb{Z}$-algèbres de prés. finie \ill{} ; $\mathrm{Mod}^{\circ}_{\mathrm{pf}}
\simeq$ schémas \struck{\ill{}} \ill{} \ill{} sur schémas affines de t.f. sur
$\mathbb{Z}$.

$M$ inv. $\Longleftrightarrow$ $M \otimes_{A} \underline{\mathrm{Hom}}_{A}(M,
A) \xrightarrow{\ \sim\ } A$ \ill{} n'est pas exprimable par $\varinjlim$ /
ne commute pas au changt de base \ill{} :
$\underline{\mathrm{Hom}}(X, Y) \to \underline{\mathrm{Hom}}(f^{*}X, f^{*}Y)$
\ill{} \struck{\ill{}} $\underline{\mathrm{Hom}}(X, Y)$ \ill{}
\note{la ligne est écrite par-dessus une autre, biffée}

\note{encadré :}
$\exists N$, $M \otimes N \simeq A$ ; $X \mapsto M \otimes X$ ; $T_{M} T_{N}
\simeq \mathrm{id}$, $T_{N} T_{M} \simeq \mathrm{id}$, donc $T_{N}$, $T_{M}$
des équivalences ; donc \ill{} ; donc $M$ et $N$ plats \add{i.e.} ; $M =
\varinjlim M_{i}$, $M_{i} \subset M$ de t.f., $M \otimes N = \varinjlim M_{i}
\otimes N$, donc $M_{i} \otimes N \subset A$, donc \ill{} \ill{}, donc
$M_{i} \otimes N = A$ ou $A$ \ill{}, \ill{} $M_{i} = M$, donc $M$ de t.f.,
\struck{\ill{}} $N$ pareil, i.e. $M$ est \ill{}.
\struck{\ill{} \ill{}} \note{deux lignes de l'encadré sont raturées au point
d'être illisibles}

\page{54}
\note{feuillets jaunis, d'une autre main de plume : le run qui commence ici
est un relevé d'énoncés sur les catégories accessibles, avec en marge des
renvois numérotés — « 8.13.2 (+ 8.12.8.1) », « 9.22 » — qui sont ceux de
SGA 4, exposé I}

\section*{Catégories $\mathrm{Ind}_{\pi}(\mathcal{C})$ ($\mathcal{C}$
petite) et $\mathrm{Fil}^{\pi}(\mathrm{Ind}_{\pi}(\mathcal{C}))$}
\note{le titre est le sien ; entre les deux il a biffé une première forme
$\mathrm{Ind}_{\pi}(\mathcal{C})_{\pi'}$}

\underline{Th.}\ $E$ une catégorie [stable par petites $\varinjlim$ ?
\add{— en \uncertain{satisfaisant} seulement $L_{\pi}$ ?}] Conditions
équivalentes :
\begin{itemize}
  \item[a)] $E$ $\pi$-accessible
  \item[b)] $E$ équivalente à une catégorie $\mathrm{Ind}_{\pi}(\mathcal{C})$,
    $\mathcal{C}$ petite
  \item[c)] La sous-catégorie $E_{\pi}$ des $E$ formée des objets
    $\pi$-accessibles est équiv. à une petite cat. et elle est strictement
    génératrice.
\end{itemize}
\struck{\ill{} filtration cardinale}, Sous ces conditions, \struck{\ill{}}
on a
\[
  E_{\pi'} \xrightarrow{\ \approx\ } \mathrm{Ind}_{\pi}(E_{\pi})_{\pi'}
  \quad \add{= \mathrm{Ind}_{\pi}(\mathcal{C})_{\pi'}},
\]
et les $(E_{\pi'})_{\pi' \geqslant \pi}$ forment une filtration cardinale ;
\struck{\ill{}} Pour les filtrations cardinales des
$(\mathrm{Fil}^{\pi})_{\pi \geqslant \pi_{0}}$ \ill{} :
\[
  \mathrm{Fil}^{\pi'}(E) = E_{\pi'} \quad \text{pour } \pi' \text{ grand.}
\]

\underline{Cor.}\ Soient $E$ et $F$ deux catégories accessibles,
\struck{Alors} $f : E \to F$ un foncteur. \struck{Alors} Conditions
équiv. :
\begin{itemize}
  \item[a)] $f$ est \struck{$\pi$-}accessible
  \item[b)] $\exists \pi$ tel que $f \simeq \mathrm{Ind}_{\pi}(f_{\pi})$, où
    $f_{\pi} : E_{\pi} \to F_{\pi}$ un foncteur \add{[où $E \simeq
    \mathrm{Ind}_{\pi}(E_{\pi})$, $F \simeq \mathrm{Ind}_{\pi}(F_{\pi})$ pour
    $\pi \geqslant \pi_{0}$]}. \struck{Pour qu'un $\pi$ convienne, il faut
    et suffit que $\pi$ \ill{}, et que $f$ soit $\pi$-accessible et
    \ill{} $E_{\pi}$ dans $F_{\pi}$}
  \item[c)] $\exists$ une petite sous-catégorie pleine $\mathcal{C}$ de $E$,
    $i : \mathcal{C} \to E$, telle que $f$ soit dans l'image essentielle des
    $i^{F}_{!}$ : $\underline{\mathrm{Hom}}(\mathcal{C}, F) \to
    \underline{\mathrm{Hom}}(E, F)$.
\end{itemize}
\struck{b) Pour \ill{} $\pi$, $f_{i}$ est \ill{}}

\underline{Corollaire}\ \struck{Supposons $f$ $\pi$-accessible. Alors
$\exists \pi_{0}$ tel que $\pi \geqslant \pi_{0} \Longrightarrow f(E_{\pi})
\subset F_{\pi}$.} \add{$E$ $\pi$-accessible, $F$ quelc. Alors}
\[
  \underline{\mathrm{Hom}}(E, F)_{\pi} \xrightarrow{\ \sim\ }
  \underline{\mathrm{Hom}}(E_{\pi}, F) \quad \add{\text{pl. fid.}}
\]
c'est une équivalence de catégories \struck{\ill{}} \add{avec petites
$\varinjlim$} \add{— si $F$ satisfait $L_{\pi}$}.

\underline{Th}\ Soit $\mathcal{C}$ une catégorie accessible \add{\struck{\ill{}}
Alors} :
\begin{itemize}
  \item[a)] (Pour \ill{} \ill{}) les foncteurs contravariants $\mathcal{C}^{\circ}
    \to (\mathrm{Ens})$ \add{petites} qui commutent aux \ill{} \ill{} sont
    \struck{\ill{}} \ill{} ; $\mathcal{C} \to (\mathrm{Ens})$ est
    pro-représentable (resp.) \ill{}
  \item[b)] les foncteurs covariants \add{s'il est} \struck{\ill{} $\pi$ est
    accessible et} exact à gauche (resp. commutant aux petites
    $\varprojlim$) et \underline{accessible} \add{stable par petites lim}.
\end{itemize}
\marginal{8.13.2 (+ 8.12.8.1)}

\underline{Corollaire}\ Soient $f : \mathcal{C} \to \mathcal{C}'$ un
foncteur, avec $\mathcal{C}$ accessible \add{petites} :
\begin{itemize}
  \item[a)] $f$ admet un adjoint : dont ssi $f$ commute aux
    \underline{$\varprojlim$} \add{petites} ;
  \item[b)] $f$ admet un \add{pro-}adjoint \struck{\ill{}} (resp. un adjoint
    \ill{}) ssi il est exact : \ill{} commute aux \add{petites}
    $\varprojlim$ (resp.) et est \underline{accessible}.
\end{itemize}

\underline{Th.}\ a) Soit $E \to B$ une catégorie fibrée, avec $B$ équiv. à
une \struck{petite} \add{petite} catégorie, les catégories fibres
accessibles, les foncteurs changement de base accessibles. Alors
$\varprojlim E/B$ et $\underline{\mathrm{Hom}}_{B}(B, E)$ sont accessibles.
En particulier, si $\mathcal{C}$ accessible et $I$ équiv. à une petite
catégorie, $\underline{\mathrm{Hom}}(I, \mathcal{C})$ est accessible, ainsi
que la sous-catégorie pleine des foncteurs qui transforment \ill{} en iso.
\marginal{9.22}

b) Mêmes hyp. sur $E, B$. Alors $\varinjlim E/B^{\circ}$ est accessible.
\marginal{à vérifier}

\page{55}
\section*{Catégories accessibles}
\note{le titre est le sien ; en haut à droite : « Réf : SGA 4 I »}

\struck{Catégorie filtrante} \add{Ensemble ordonné} grande devant $\pi$.
Conditions $L_{\pi}$, $L$ \\
Foncteur $\pi$-accessible, accessible \\
Objet $\pi$-accessible, objet accessible d'une catégorie.
\marginal{9.1 : 9.3 — en regard des trois lignes}

Catégorie $\pi$-accessible : Satisfait $L_{\pi}$, et admet une \add{petite}
sous-catégorie génératrice \add{\underline{str.}} \struck{telle que} dont
les objets sont $\pi$-accessibles ; accessible : $\exists \pi$ tel que
$\mathcal{C}$ $\pi$-accessible

\marginal{en biais : (vérifier 9.4 ?) — c'est la bonne notion si $E$ stable
par petites $\varinjlim$ ; sinon sont les bonnes catégories les
$\mathrm{Ind}_{\pi}(\mathcal{C})$ avec $\mathcal{C}$ petite \ldots — 9.6}

\underline{Prop}\ \struck{\ill{}} $\underline{\mathrm{Hom}}(E, F)_{\pi}$,
$F$ satisfaisant $L_{\pi}$
\begin{itemize}
  \item[a)] stable par \add{petites} $\varinjlim$ qui sont représentables
    dans $F$
  \item[b)] Si $F$ $\pi$-accessible (\ill{}) alors pour tt catégorie $I$
    \add{avec card $I \leqslant \pi$} et les lim de type $I$ rep. dans $F$,
    $\underline{\mathrm{Hom}}(\ )_{\pi}$ stable par lim type $I$.
\end{itemize}
\struck{\ill{}}

\underline{Cor 1}\ Cas accessible

\underline{Cor 2}\ \struck{\ill{}} Si $F$ est (resp) \struck{\ill{}}
\add{$\pi$-}accessible, \struck{\ill{}} pour tt $\pi' \geqslant \pi$,
$E_{\pi'}$ est stable par lim de \struck{type} $I$, avec card $I \leqslant
\pi'$. [Les objets \ill{} stables par \add{petites} $\varinjlim$]

\underline{Cor 3}\ Si \struck{accessible} $E$ est accessible \add{$F$
stable}, tout objet de $E$ est accessible.

Réciproque partielle : \underline{Prop}\ \struck{\ill{}} $E$ catégorie
satisfaisant la condition $L$ et : \struck{a) $E$ admet une petite
sous-catégorie str. gén.}
\begin{itemize}
  \item[a)] $E$ stable par noyaux de doubles flèches, et $\mathrm{Ker}$ est
    accessible (p. ex. commute aux \add{petites} $\varinjlim$ filtrantes)
  \item[b)] $E$ stable par produits fibrés \marginal{($E$ admet une pet.
    ss-cat. str. génératrice, et)}
  \item[c)] [$E$ admet une petite sous-catégorie \struck{génératrice}
    génératrice par épim. stricts universels (p. ex. \struck{\ill{}} les
    familles épimorphiques strictes sont épimorphiques strictes
    \struck{\ill{}} universelles ; i.e. si les familles épimorphiques
    \struck{\ill{}} strictes \ill{}]
    \note{le c) est repris trois fois, la formule finale surchargeant les
    deux premières}
\end{itemize}
Alors $E$ est accessible (i.e., en présence de a)) tout objet de $E$ est
accessible

\struck{\ill{}}

\underline{Cor}\ $E$ catégorie stable par petites $\varinjlim$, \struck{et
telles que les familles épimorphiques strictes universelles} Conditions
équivalentes
\begin{itemize}
  \item[a)] $E$ est accessible
  \item[b)] $E$ admet une petite sous-catégorie \struck{génératrice}
    \add{par épim. str. universels} \add{comme} génératrice
    \struck{\ill{}} (\add{petites} lim) et les foncteurs Ker \ill{} sur la
    catégorie des doubles flèches des $E$ est accessible.
\end{itemize}
[\textbf{N.B.}\ $E$ accessible $\Longrightarrow$
$\underline{\mathrm{Hom}}(\mathcal{C}, E)$ accessible pour tte petite
catégorie $\mathcal{C}$] ; $\exists$ ($E$ $\pi$-accessible) $E_{\pi}$ petite
sous-catégorie gén. par épim. stricts universels]

\underline{Ex.}\ Dans les topos, \struck{\ill{}} \add{à $\varinjlim$
filtrantes exactes et à petites} \ill{} des catégories accessibles ; une
catégorie \struck{\ill{}} : petite famille génératrice \struck{\ill{}}, est
accessible. Mais $(\mathrm{Ens})^{\circ}$ et $(\mathrm{Ab})^{\circ}$ ne sont
pas accessibles, p. ex. \add{$\mathrm{Ab}^{\circ}$}, \add{p. ex.} par que
dans ces catégories \struck{\ill{}} \struck{il n'y a pas de sous-catégorie
génératrice par épim. stricts universels} \ill{} n'est pas accessible [i.e.
les foncteurs Ker dans Ens, (Ab) ne commutent pas aux \ill{} par
\uncertain{lim} : une $\varinjlim$ grande devant $\pi$ d'épim n'est pas
néc. un épim.]

\underline{Filtrations cardinales} d'une catégorie $E$ \struck{\ill{}} p. ex.
$\mathrm{Fil}^{\pi}(E)$, $\pi \geqslant \pi_{0}$ :
\begin{itemize}
  \item[a)] Les $\mathrm{Fil}^{\pi}(E)$ équivalentes à petite catégorie
  \item[b)] \struck{Fil} $E$ satisfait $L_{\pi_{0}}$, $\mathrm{Fil}^{\pi}$
    stable par $\varinjlim$ filtrantes grandes avec $I$ grand devant
    $\pi_{0}$ et card $I \leqslant \pi$
  \item[c)] Si $\pi' \geqslant \pi \geqslant \pi_{0}$, $H$ élément $X$ de
    $\mathrm{Fil}^{\pi'}$ est de la forme $\varinjlim_{i \in I} X_{i}$, $X_{i}
    \in \mathrm{Fil}^{\pi}$, $I$ grand devant $\pi$, card $I \leqslant \pi'$
\end{itemize}
\marginal{9.12 (condition c) renforcée)}

\page{56}
\begin{tikzcd}
\mathcal{C} \arrow[r, "\varphi"] & E \arrow[r, "\hat{\varphi}"] \arrow[dr, dashed] & \hat{\mathcal{C}} \\
& & \hat{E} \arrow[u, dashed, "\hat{\varphi}^{*}"']
\end{tikzcd}

\[
  \mathrm{Hom}_{E}(\varphi(Y), X) \simeq \mathrm{Hom}_{\hat{\mathcal{C}}}
  (Y, \hat{\varphi}(X))
\]
\struck{\ill{}} \note{une seconde formule, biffée au point d'être
illisible}

Si ds $E$ $\exists$ limites inductives \add{(petites)}, on a aussi :
\[
  \bar{\varphi} : \hat{\mathcal{C}} \longrightarrow E, \qquad
  \bar{\varphi}(F) = \varinjlim_{\mathcal{C}_{/F}} \varphi(X)
\]
le choix de $\varphi$ équivaut à celle de $\bar{\varphi}$ commutant aux
petites $\varinjlim$
\[
  \mathrm{Hom}_{E}(\bar{\varphi}(F), X) \simeq
  \mathrm{Hom}_{\hat{\mathcal{C}}}(F, \hat{\varphi}(X))
\]

\underline{Conditions équivalentes}
\begin{itemize}
  \item[a)] $\forall X \in E$, $\varinjlim_{Y \in \mathcal{C}_{/X}}
    \varphi(Y) \xrightarrow{\ \sim\ } X$
  \item[b)] $\hat{\varphi} : E \to \hat{\mathcal{C}}$ — pl. fidèle
\end{itemize}
\marginal{en accolade à droite de a) et b) : On dit alors que $\varphi$ est
strictement générateur.}
\begin{itemize}
  \item[b')] (Si $E$ stable par \add{petites} $\varinjlim$) $\hat{\mathcal{C}}
    \to E$ un foncteur de localisation [\add{on \uncertain{sait} : un.
    surj.}]
  \item[c')] \struck{\ill{}} (Si $E$ stable par \add{petites} $\varinjlim$
    filtrantes, et les $\mathcal{C}_{/X}$ ($X \in \mathrm{Ob}\,E$) filtrant,
    p. ex. $\mathcal{C}$ stable par \underline{lim} finies et $\varphi$
    exact à g.) : $\tilde{\varphi} : \mathrm{Ind}(\mathcal{C}) \to E$ induit
    par $\bar{\varphi}$ est un foncteur de localisation
  \item[c)] (Si $\varphi$ admet un ind-adjoint, i.e. $\hat{\varphi}$ prend
    ses valeurs dans $\mathrm{Ind}(\mathcal{C})$) ${}^{i}\varphi : E \to
    \mathrm{Ind}(\mathcal{C})$ est pl. fidèle.
  \item[d)] (Si $\varphi$ admet un adjoint : \ill{} $\psi$, i.e.
    $\hat{\varphi}$ prend ses valeurs dans $\mathcal{C}$) $\psi$ est pl.
    fidèle
  \item[d')] \struck{$\varphi$ est un \ill{}} (\add{$\varphi$ est}
    (\ill{} conditions) $\varphi$ est un foncteur de localisation.
  \item[e)] (Si $\varphi$ pl. fidèle) $\varphi(\mathcal{C})$ est génératrice
    dans $E$ par épim. stricts.
  \item[f)] Pour toute catégorie $F$ stable par petites $\varinjlim$,
    désignant par $\underline{\mathrm{Hom}}'(E, F)$ la sous-catégorie pleine
    de $\underline{\mathrm{Hom}}$ formée des foncteurs qui commutent aux
    petites $\varinjlim$, le foncteur restriction
    \[
      \beta_{F} : \underline{\mathrm{Hom}}'(E, F) \longrightarrow
      \underline{\mathrm{Hom}}(\mathcal{C}, F)
    \]
    est pl. fidèle
  \item[g')] \struck{Avec les notations précédentes, \ill{}}
  \item[g)] (Si $\mathcal{C}$ stable par $\varinjlim$ finies, et $\varphi$
    y commute \add{et est pl. fid.}) $\beta_{F}$ précédent induit une
    équivalence entre la catégorie $\underline{\mathrm{Hom}}'(E, F)$ et la
    sous-catégorie pleine de $\underline{\mathrm{Hom}}(\mathcal{C}, F)$
    formée des $f : \mathcal{C} \to F$ \ill{} qui commutent aux \ill{}
    \add{$\varinjlim$} finies : \ill{} dont \ill{} induisent \ill{} dans
    $F$ \ill{}
  \item[] \struck{Si $\mathcal{C}$ stable par lim finies, $\mathcal{C}$
    \ill{} des petites lim, les \ill{} filtrantes sont exactes, $\varphi$
    est exact \ill{}}
  \item[h)] \struck{\ill{}} (Si $E$ un topos, $\mathcal{C}$ stable par
    $\varprojlim$ finies, $\varphi$ exact à g.) $\varphi$ définit $E$ comme
    un \underline{sous-topos} de $\hat{\mathcal{C}}$.
  \item[] \struck{i) \ill{}}
\end{itemize}

\page{58}
$E$ \underline{catégorie}. Conditions équivalentes
\begin{itemize}
  \item[a)] \struck{Il existe une} $E$ est stable par petites $\varinjlim$ et
    admet une petite \add{sous-}catégorie pleine str. génératrice
  \item[a')] $E$ est stable par petites $\varinjlim$, et il existe une petite
    catégorie $\mathcal{C}$ et un foncteur $\varphi : \mathcal{C} \to E$
    str. générateur
  \item[a'')] Il existe une petite catégorie $\mathcal{C}$, et un foncteur
    pleinement fidèle $E \to \hat{\mathcal{C}}$ qui admet un adjoint à gauche
    $\hat{\mathcal{C}} \to E$.
\end{itemize}
[Dans une telle catégorie, tt foncteur \add{contravariant} (à valeurs dans
(Ens)) qui commute aux petites $\varinjlim$ est représentable, donc il existe
dans $E$ des petites $\varprojlim$, et dans a''), le foncteur $E \to
\hat{\mathcal{C}}$ y commute, i.e. $E$ s'identifie à une sous-catégorie
stable par $\varprojlim$ de $\hat{\mathcal{C}}$. Si $E$ est accessible [ce
qui est plus fort que les conditions équivalentes a) \ldots], alors
\struck{\ill{}} un foncteur $E \to \hat{\mathcal{C}}$ \add{\ill{}} un adjoint
à g. ssi il est \underline{accessible} et commute aux \add{petites}
$\varprojlim$]

\underline{Question}\ Une catégorie satisfaisant les conditions a) est-elle
\underline{accessible}, i.e. ses éléments sont-ils accessibles ? Comme, sous
les conditions de a''), $\hat{\mathcal{C}}$ est accessible, on voit qu'il
\add{faut et} suffit que le foncteur $E \to \hat{\mathcal{C}}$ soit
accessible, ou ce qui revient au même, qu'il existe un cardinal $\pi$ tel que
$E$ soit stable dans $\hat{\mathcal{C}}$ par $\varinjlim$ filtrantes grandes
devant $\pi$.

\marginal{en biais dans la marge gauche : Contre-exemple : $(\mathrm{Ens})^{\circ}$,
elle \ill{} à des limites inductives strictes \ill{} épimorphismes
\uncertain{stricts} \ill{} (Ex $\geqslant 2$) \ill{} les \ill{} de \ill{} ne
sont pas accessibles \ill{} les \uncertain{produits} !}

\underline{Exemples}
\begin{itemize}
  \item[a)] Topos \struck{\ill{}} \add{a') opposé de topos \struck{\ill{}}}
  \item[b)] $T(\mathrm{Ens})$, $T$ une structure algébrique définissable par
    \underline{petites} $\varprojlim$
  \item[c)] $T(\mathcal{C})$, où $\mathbf{T}$ est comme dans b), et
    $\mathcal{C}$ une catégorie \underline{accessible}
\end{itemize}
\struck{\ill{}}

Conséquence : 1) Si $T \to T'$ est un hom. de structures algébriques, d'où
$T(\mathcal{C}) \to T'(\mathcal{C})$ foncteur \underline{accessible} entre
catégories accessibles, commutant aux $\varprojlim$ petites, donc $\exists$
adjoint à gauche (structures $T'$-$T$-libres) ; 2) Si $\mathcal{C} \to
\mathcal{C}'$ foncteur accessible commutant aux pet. $\varprojlim$ entre cat.
acc., $T(\mathcal{C}) \to T(\mathcal{C}')$ est \underline{accessible}
commutant aux petites $\varprojlim$, donc $\exists$ adjoint à g. \ldots

\[
  \underline{\mathrm{Hom}}^{*}(\mathcal{C}^{\circ} \times \mathcal{C}'^{\circ},
  \mathrm{Ens})
\]
À prouver : Si $\mathcal{C}$, $\mathcal{C}'$ accessibles, alors la
ss-catégorie \add{des} $\underline{\mathrm{Hom}}(\mathcal{C}^{\circ} \times
\mathcal{C}'^{\circ}, \mathrm{Ens})$ formée des bifoncteurs qui commutent aux
$\varprojlim$ \add{petites} en chaque argument, ($\simeq
\underline{\mathrm{Hom}}'(\mathcal{C}^{\circ}, \mathcal{C}') \simeq
\underline{\mathrm{Hom}}'(\mathcal{C}'^{\circ}, \mathcal{C})$, où
$\underline{\mathrm{Hom}}'$ désigne les foncteurs qui commutent aux petites
$\varprojlim$) est accessible, et des foncteurs commutant aux $\varprojlim$
$\mathcal{C}_{1} \to \mathcal{C}$, $\mathcal{C}'_{1} \to \mathcal{C}'$
définissent un foncteur accessible entre les
$\underline{\mathrm{Hom}}^{*}$. Or il est assez clair que plus
\uncertain{généralement}, si $\mathcal{D}$ \add{variable} \struck{stable par}
\ill{} (et pour les foncteurs (qui) commutent \ill{})
\[
  \underline{\mathrm{Hom}}^{*}(\mathcal{C}^{\circ} \times \mathcal{C}'^{\circ},
  \mathcal{D}) \simeq T_{\mathcal{C}^{\circ}} T_{\mathcal{C}'^{\circ}}(\mathcal{D})
  \simeq T_{\mathcal{C}'^{\circ}} T_{\mathcal{C}^{\circ}}(\mathcal{D})
\]
est une « théorie algébrique » que, définie par $(\mathcal{C} \times
\mathcal{C}')^{\circ}$ [$\mathcal{C} \times \mathcal{C}'$ universel pour
« couplages » $\mathcal{C} \times \mathcal{C}' \to \mathcal{C}''$ commutant
aux $\varprojlim$ en chaque argument \ldots]

\page{59}
\underline{Lemme}\ $\mathcal{C}$ une petite catégorie, $\pi$ cardinal, $E =
\mathrm{Ind}_{\pi}(\mathcal{C})$ la sous-catégorie pleine de
$\underline{\mathrm{Ind}}(\mathcal{C})$ formée des repl. ind. \add{indexés par
des ens. ordonnés \add{petits}} grands devant $\pi$. Alors (c'est une
$\mathcal{U}$-catégorie \ill{})
\begin{itemize}
  \item[(i)] $E$ satisfait $L_{\pi}$, (et l'inclusion $E \to
    \mathrm{Ind}(\mathcal{C})$ \struck{commute} est $\pi$-accessible.
  \item[(ii)] Les objets de $E$ qui sont $\pi$-accessibles sont ceux qui sont
    isom. à des facteurs directs d'objets de $\mathcal{C}$ (donc forment une
    petite catégorie équivalente \ill{} l'enveloppe Karoubienne de \ill{} de
    $E$, compte tenu que $\mathrm{Ind}_{\pi}(\mathcal{C})$ est stable par
    facteurs directs, cf (i)) ; En particulier, si $\mathcal{C}$ est
    karoubienne, \struck{i.e. \ill{} « les pts isolés » de $E$ \ill{}}
    $\mathcal{C}$ \add{dans} $E_{\pi}$ \add{des} objets de $\mathcal{C}$.
  \item[(iii)] $\mathcal{C}$ est strictement génératrice dans $E$
  \item[(iv)] Pour tt objet $X$ de $E$, $\mathcal{C}_{/X}$ est une catégorie
    filtrante \add{(petite)} grande devant $\pi$
\end{itemize}
\struck{\ill{}}

\underline{Réciproquement} : Soit $E$ une \struck{\ill{}} catégorie \add{$\pi$
un cardinal \add{petit}}, \struck{la forme $\mathrm{Ind}(\mathcal{C})$,
\add{$\mathcal{C}$ petite}, \ill{} et satisfaisant} \add{et une
$\mathcal{U}$-catégorie et} — Alors $E$ est $\approx$ à une catégorie du
\add{\ill{}} :
\begin{itemize}
  \item[(i)] $E$ satisfait $L_{\pi}$
  \item[(ii)] La sous-catégorie \add{$E_{\pi}$} de $E$ formée des objets
    $\pi$-accessibles est \struck{petite} petite, et strictement
    génératrice
  \item[(iii)] Pour tt objet $X$ de $E$, la catégorie $E_{\pi/X}$ est filtrante
    et grande devant $\pi$.
\end{itemize}
\underline{N.B.}\ La condition (iii) est \struck{vérifiée} \add{conséquence de
(i), (ii)} (si $E$ stable par \add{petites} $\varinjlim$ (\ill{} $E_{\pi}$
est stable par $\varinjlim$ de cardinal $\leqslant \pi$)

\struck{Soit $\pi'$ un cardinal, et}

Reprenons $\mathrm{Ind}_{\pi}(\mathcal{C})$, soit $\pi'$ un cardinal
$\geqslant \pi$, et soit $\mathrm{Ind}_{\pi}(\mathcal{C})_{\pi'}$ \add{$=
\mathrm{Fil}^{\pi'}(E)$} la sous-catégorie str. pleine des \ill{} de
$\mathrm{Ind}_{\pi}(\mathcal{C})$ formée des ind-objets indexés par des $I$
tels que a) $I$ grand devant $\pi$ b) \struck{card Ob $I$} \add{Ob $I$
\ill{}} ensemble cofinal de cardinal $\leqslant \pi'$. Alors
\begin{itemize}
  \item[a)] Les objets de $\mathrm{Ind}_{\pi}(\mathcal{C})_{\pi'}$ \add{$=
    \mathrm{Fil}^{\pi'}(E)$} dans $E$ sont $\pi'$-accessibles \add{(en
    objets)} : $\mathrm{Fil}^{\pi'}(E) \subset E_{\pi'}$
    \marginal{on \ill{} de $E$ (tt objet de $E$ est accessible)}
  \item[b)] Si $\pi'$ est de la forme $2^{c}$ ($c \geqslant \pi$), alors la
    réciproque est vraie i.e. $\mathrm{Fil}^{\pi'}(E) = E_{\pi'}$
  \item[c)] La filtration de $E$ par les \struck{$\mathrm{Ind}$}
    $\mathrm{Fil}^{\pi'}(E)$ ($\pi' \geqslant \pi$) \add{est une}
    « filtration cardinale » \ill{} dans \ill{} (Si $E$ est stable par
    \struck{petites} $\varinjlim$) i.e. \struck{$E$ stable par $\varinjlim$
    de card $\leqslant \pi'$} \add{$E$ stable par lim de card $\leqslant
    \pi'$} (\ill{} \ill{} $\pi' \geqslant \pi$)
  \item[d)] \struck{Pour tt cardinal $\pi' \geqslant \pi$} \add{(Si $E$
    stable par $\varinjlim$)} \struck{Pour tt cardinal $\pi' \geqslant \pi$,
    \ill{} $E \approx \mathrm{Ind}(E_{\pi'})$} \add{$(E_{\pi'})_{\pi'
    \geqslant \pi}$ forment une} $E$ est équivalente à
    $\mathrm{Ind}(E_{\pi'})$. \struck{\ill{}} \add{« filtration
    cardinale »}
\end{itemize}
\note{les points c) et d) sont surchargés de renvois et de traits ; l'ordre
ci-dessus suit les flèches}
\marginal{Pourtant, on a $\mathrm{Fil}^{\pi'}(E) \neq E_{\pi'}$ — et,
encerclé : prouver c) avant b)}
\marginal{encadré : cf SGA 4 I 9.66 pour $\supset$ — et : résulte de la
réciproque plus haut ; on a égalité si $E$ stable par $\varinjlim$}

\page{60}
\underline{Axiome c) des filtrations cardinales}

\underline{Pour $(\mathrm{Fil}^{\pi})_{\pi'}$} :
\[
  X = \varinjlim_{I} X_{i}, \qquad X_{i} \in \mathcal{C},\ I \text{ grand
  devant } \pi. \quad \struck{\ill{}}
\]
Soit $I'$ l'une des \struck{sous-catégories pleines de $I$} \add{ordonnés par
l'inclusion} \add{filtrantes de $I$ \struck{grandes devant $\pi$} (telles
que} card(Ob $I'$) $\leqslant \pi'$, \struck{\ill{}} Pour tt tel $I'$, soit
$X_{I'} = \varinjlim_{i \in I'} X_{i}$ (limite inductive dans $E =
\mathrm{Ind}_{\pi}(\mathcal{C}) \subset \mathrm{Ind}(\mathcal{C})$)
\struck{\ill{}}. Je dis que $J$ est filtrant et grand devant $\pi'$. En
effet, si $(I_{\alpha})_{\alpha \in \Lambda}$ est une famille de
sous-catégories de $J$ \struck{grandes devant $\pi$}, de cardinal $\leqslant
\pi'$, avec card $\Lambda \leqslant \pi'$, la réunion \add{\ill{}} des
$I_{\alpha}$ est de cardinal $\leqslant \pi'$, et \struck{\ill{}} est \ill{}
dans \ill{} \ill{} de cardinal $\leqslant \pi'$ \struck{\ill{}} \struck{\ill{}
grand devant $\pi$, pour toute partie de $I'$ de \ill{}} les $X_{I'}$ dans
$\mathrm{Fil}^{\pi'}$,
\[
  X = \varinjlim_{I' \in J} X_{I'} .
\]
Enfin, si $X \in \mathrm{Fil}^{\pi''}$, on peut prendre $I$ tel que card $I
\leqslant \pi''$, donc Card $J \leqslant \pi''^{\pi'}$.

\marginal{en regard, colonne de droite, sous « $\forall \pi' \geqslant \pi$ » :
Pour $X \in E$ de la forme $\varinjlim_{J} X_{j}$, $X_{j} \in
\mathrm{Fil}^{\pi'}$, $J$ grand devant $\pi'$, et si $X \in
\mathrm{Fil}^{\pi''}$, on peut prendre card $J \leqslant \pi''^{\pi'}$
($= \pi''$ ?) (\underline{N.B.}\ si $\pi'' \geqslant 2^{\pi'}$, \ldots)}

\underline{Axiome des filtrations cardinales pour les $(E_{\pi'})$}. On a vu
que, si $X \in \mathrm{Ob}\,E_{\pi''}$, on \add{sans diff. \ill{}}
\struck{\ill{} $\pi'' \geqslant \pi'$} \add{(SGA 4 I 9.16)} On a
\[
  X = \varinjlim_{J} X_{j}, \qquad X_{j} \in \mathrm{Fil}^{\pi'}(X) \cap
  E_{\pi'},\ J \text{ grand devant } \pi',
\]
\add{à prouver que $\{\,X_{j} \in \mathrm{Ob}\,E_{\pi'}$, $J$ grand devant
$\pi'$, card $J \leqslant \pi''^{\pi'}\,\}$} donc tout revient à prouver
\uncertain{que} card $J \leqslant \pi''^{\pi'}$. \struck{Ceci \ill{} $\pi''
\geqslant \pi'$, sauf \ill{} $\leqslant \mathrm{Fil}$}
\note{le paragraphe entier est barré de traits diagonaux}

\struck{\underline{Cor} Pour \add{supposons que}} \note{amorce encerclée et
abandonnée}

\underline{Cor}\ Soit $E$ une $\mathcal{U}$-cat. stable par petites
$\varinjlim$ \add{\ill{}}. Pour \add{que} $E$ soit équivalente à
\struck{\ill{}} $\mathrm{Ind}_{\pi}(\mathcal{C})$, avec $\mathcal{C}$ petite,
\add{(mieux : $E$ \ill{} par $\varinjlim$ de card $\leqslant \pi$)} il f. et
s. que $E$ admette une petite catégorie strictement génératrice formée
d'objets accessibles. \marginal{encerclé : qui il en est $\pi$ ? il tel
\ill{}} Ce sont ces $E$ qu'il faudrait appeler « accessibles ».

\end{document}
